% GENERATED FILE. Edit canonical lessons, then run npm run generate:books.
% canonical-source-hash: fnv1a64:771849e1943269f5
% canonical-lessons: TA-C27-maalai

\chapter{Evening}
\label{ch:ta-evening}

This chapter is generated from the canonical micro-lessons used by Language Ladder.
Each section stays independently resumable and preserves the authored prerequisite order.

\section[\ta{மாலை}]{\ta{மாலை} (mālai) — "evening," and a garland that only looks like the same word}

\label{lesson:TA-C27-maalai}



\begin{quote}
\textbf{Warm-up.} {[}PAUSE 2s{]} Last lesson's morning word had an uncertain origin. This evening word is different — confidently native — but it comes with its own honest twist: it looks exactly like a completely unrelated word.
\end{quote}

\begin{cousinweb}
\textbf{\ta{மாலை}} (\textbf{mālai}) — "\textbf{evening}" — happens to be spelled identically to another, completely \textbf{unrelated} Tamil word: \textbf{\ta{மாலை}} meaning "\textbf{garland, wreath, necklace}." Wiktionary treats these as \textbf{two separate etymologies}, not one word with two senses. The garland sense traces to Old Tamil and is compared to Sanskrit \textbf{\dv{माला}} (\emph{mālā}) — with a genuinely interesting note: it's "\textbf{probably borrowed FROM Dravidian languages}" \textbf{into} Sanskrit, not the other way around. A rare reverse-direction loan, worth noticing against this whole project's usual Sanskrit-into-Dravidian pattern.
\end{cousinweb}

\begin{cousinweb}
The \textbf{evening} sense of \textbf{\ta{மாலை}} is a \textbf{separate} Old Tamil word entirely, compared instead to \textbf{\ta{மால்}} (\emph{māl}) — which Wiktionary traces to \textbf{Proto-Dravidian} \emph{\textbf{*mā/*māl}}, with senses including "\textbf{\ta{கருமை}}" ("\textbf{blackness, darkness}"), plus archaic senses of "illusion, confusion" and, as a proper noun, the Hindu deity \textbf{Vishnu}. Be honest about the exact mechanism here: Wiktionary's own etymology for "evening" only says "\textbf{compare māl}" — it does \textbf{not} explicitly state that "evening" derives specifically from the \textbf{darkness} sense rather than some other sense of this multi-meaning root. The connection (evening = the time it grows dark) is natural and plausible, but not nailed down with full certainty by the source itself. What IS certain: this is confidently \textbf{native} Dravidian vocabulary, a genuine contrast with TA-C26's uncertain \ta{காலை}.
\end{cousinweb}

\subsection*{Guided Practice}
{[}PAUSE 1s{]}

\begin{itemize}
\raggedright
  \item {[}YOU SAY: "mālai" — "evening," confidently native Dravidian vocabulary{]}
  \item {[}YOU SAY: the homonym twist — a completely separate word, mālai, means "garland," probably borrowed INTO Sanskrit from Dravidian{]}
  \item {[}YOU SAY: the honest hedge — "evening" plausibly connects to māl's "darkness" sense, but Wiktionary doesn't spell out that exact link{]}
\end{itemize}

\begin{tcolorbox}[breakable,colback=teal!4,colframe=teal!35!black,title={Wrap-up recall}]
{[}PAUSE 3s{]} Are evening-mālai and garland-mālai the same word with two meanings? (\textbf{No} — Wiktionary treats them as two entirely separate etymologies; a genuine homonym, not one word stretched.) Which direction did the garland sense's Sanskrit connection likely travel? (\textbf{Dravidian into Sanskrit} — a reverse-direction loan, per Wiktionary's own "probably borrowed from Dravidian languages" note.) Is it fully certain that "evening" derives from māl's "darkness" sense specifically? (\textbf{No} — plausible and natural, but Wiktionary only says "compare māl," without spelling out that exact mechanism.)
\end{tcolorbox}
